\section{The complete cyclotomic lattice}
Let $p$ be an odd prime, let $\zeta$ be a primitive $p$th root of unity,
and put $O=\mathbb Z[\zeta]$ and $\pi=\zeta-1$.
% BEGIN UF.Lattice
\begin{theorem}[Cyclotomic coordinates]\label{UF.Lattice}
For every $n\geq1$, the ring $A_n=O/\pi^nO$ has $p^n$ elements.
Write uniquely
\[
 n=(e-1)(p-1)+r,\qquad e\geq1,\quad 1\leq r\leq p-1.
\]
As an additive group with the indicated ordered generators,
\begin{equation}\label{eq:coordinates}
 A_n=\gen{1}\oplus\gen{\pi}\oplus\cdots\oplus\gen{\pi^{p-2}}
 \cong (\mathbb Z/p^e\mathbb Z)^r
 \oplus(\mathbb Z/p^{e-1}\mathbb Z)^{p-1-r}.
\end{equation}
A modulus-one factor is trivial. Complex conjugation induces an
involutory automorphism of $A_n$. Multiplication by $\zeta$ has order
$p$ when $n\geq2$, and for $j\not\equiv0\pmod p$ it satisfies
$\sum_{i=0}^{p-1}\zeta^{ij}=0$.
\end{theorem}
\begin{proof}
The polynomial
\[
 \Phi_p(1+X)=p+\binom p2X+\cdots+
 \binom p{p-1}X^{p-2}+X^{p-1}
\]
is Eisenstein at $p$. Thus $O$ is an integral domain with
$\mathbb Z$-basis $1,\pi,\ldots,\pi^{p-2}$. Substitution
$\zeta=1$ gives $O/\pi O\cong\mathbb F_p$.
Multiplication by $\pi^j$ induces an isomorphism
$O/\pi O\longrightarrow\pi^jO/\pi^{j+1}O$:
surjectivity is immediate and injectivity follows by cancellation
in the domain. The successive indices are all $p$, whence
\begin{equation}\label{eq:latticeorder}
 |O/\pi^nO|=p^n.
\end{equation}

The identity $p=\prod_{j=1}^{p-1}(1-\zeta^j)$ gives
$pO=\pi^{p-1}O$. Here equality, rather than inclusion alone, follows
because each ratio $(1-\zeta^j)/(1-\zeta)$ is a unit of $O$.
It belongs to $O$ by a geometric sum. The automorphism
$\zeta\mapsto\zeta^j$ identifies the norms of numerator and
denominator, both equal to $p$, so the ratio has norm one.
Its multiplication matrix on the integral basis has determinant one;
the adjugate gives an integral inverse and hence a unit of $O$.

For $1\leq r\leq p-1$ define the integral lattice
\[
 L_r=p\mathbb Z\,1+\cdots+p\mathbb Z\,\pi^{r-1}
       +\mathbb Z\,\pi^r+\cdots+\mathbb Z\,\pi^{p-2}.
\]
Every displayed generator is in $\pi^rO$, using
$pO=\pi^{p-1}O$ for the first group of generators.
Both $L_r$ and $\pi^rO$ have index $p^r$ in $O$, the latter by
\eqref{eq:latticeorder}; thus $L_r=\pi^rO$.
Now $\pi^nO=p^{e-1}\pi^rO$, which gives every coordinate modulus
in \eqref{eq:coordinates}. In particular it gives the individual
orders, including trivial listed generators when $e=1$.

The equation $\overline\pi=-\zeta^{-1}\pi$ shows that conjugation
preserves $\pi^nO$. Multiplication by $\zeta$ has $p$th power the
identity. For $n\geq2$ it is not the identity, since its difference
from the identity sends $1$ to $\pi\notin\pi^nO$; the latter follows
from the strict successive indices above. Its order is therefore $p$.
Finally $1+\zeta+\cdots+\zeta^{p-1}=0$; multiplication of exponents
by a nonzero residue $j$ permutes these terms.
\end{proof}
% END UF.Lattice

% BEGIN UF.Conjugation
\begin{corollary}[Conjugation and displacement]\label{UF.Conjugation}
Let $A_n$ be written additively and let an element $\beta$ of order
$p$ act by $v^\beta=\zeta v$. Put $s_i=\pi^{i-1}$ for $i\geq1$.
Then $[s_i,\beta]=s_{i+1}$ in $A_n\rtimes\gen{\beta}$, and
\begin{align}
 \overline{s_i}&=(-1)^{i-1}\zeta^{1-i}s_i,
       &&i\geq1,\label{eq:conjugation}\\
 \zeta^{1-i}s_i&=\sum_{j=i}^{p+1}
       \binom{p+1-i}{j-i}s_j,
       &&1\leq i\leq p.\label{eq:binomial}
\end{align}
The first expression on the left of the commutator identity is group
notation; the two displayed equations are additive ring identities.
\end{corollary}
\begin{proof}
The commutator is $(\zeta-1)s_i=\pi s_i$.
Equation~\eqref{eq:conjugation} follows by raising
$\overline\pi=-\zeta^{-1}\pi$ to the power $i-1$.
For \eqref{eq:binomial}, use $\zeta^{1-i}=\zeta^{p+1-i}$ and
expand $(1+\pi)^{p+1-i}\pi^{i-1}$. No independence of the
additional $s_p,s_{p+1}$ is asserted or needed.
\end{proof}
% END UF.Conjugation
